% Compilar con XeLaTeX o LuaLaTeX (en Overleaf: Menú > Compilador).
\documentclass[12pt,letterpaper]{article}

\usepackage{iftex}
\ifPDFTeX
  \usepackage[utf8]{inputenc}
  \usepackage[T1]{fontenc}
  \usepackage{lmodern}
  \usepackage[spanish]{babel}
\else
  \usepackage{fontspec}
  \usepackage{polyglossia}
  \setmainlanguage{spanish}
\fi

\usepackage[margin=2.5cm]{geometry}
\usepackage{amsmath}

\pagestyle{empty}

\title{Reto III.3: asistencia gravitacional en Júpiter\\ Identificación del problema}
\author{Frankly Giovanni Poveda Pinzón \and Camila Castillo López}
\date{2 de octubre de 2026}

\begin{document}
\maketitle
\thispagestyle{empty}

\section{El problema}

Se quiere diseñar un sobrevuelo de Júpiter que cambie la órbita de una sonda que viene de
la Tierra para que pueda continuar hasta Saturno. El problema se debe resolver con dos
aproximaciones: cónicas empalmadas, donde en cada tramo actúa un solo cuerpo, e integración
numérica con todos los cuerpos actuando a la vez. La pregunta que pide el reto es qué
cantidad se conserva durante el encuentro en cada sistema de referencia y cuánto error
introduce la aproximación de cónicas empalmadas.

\section{El sistema}

Los cuerpos son el Sol, Júpiter, Saturno y la sonda. La masa de la sonda es despreciable,
así que no afecta a los planetas y se trata de un problema restringido: los planetas siguen
órbitas conocidas y lo que se busca es el movimiento de la sonda. Si se trabaja en el plano
de la eclíptica quedan dos grados de libertad, la posición $(x,y)$ de la sonda.

Las escalas son muy distintas. Júpiter está a unas 5.2~UA del Sol y se mueve a unos
13~km/s. La región donde domina su gravedad (esfera de influencia, $r \approx a_J
(m_J/M_\odot)^{2/5}$) es de unos 48 millones de km, que son unos 670 radios de Júpiter y
menos del 7\,\% de su distancia al Sol. Una sonda que llega con unos 8~km/s relativos a
Júpiter tarda del orden de 140 días en cruzar esa región, mientras que el viaje completo
dura años. Esta separación de escalas es lo que, en principio, justifica empalmar cónicas.

\section{Qué se conserva en cada sistema de referencia}

En el sistema de Júpiter, si se ignora el Sol durante el encuentro, la sonda describe una
hipérbola. Se conservan la energía y el momento angular respecto a Júpiter, por lo que la
rapidez de llegada y la de salida lejos del planeta son iguales ($v_\infty$). Lo único que
cambia es la dirección, que gira un ángulo $\delta$ dado por
\[
e = 1 + \frac{r_p v_\infty^2}{\mu_J}, \qquad \sin\frac{\delta}{2} = \frac{1}{e},
\]
donde $r_p$ es la distancia de máximo acercamiento.

En el sistema del Sol la velocidad de la sonda es $\vec v = \vec v_J + \vec v_\infty$. Al
girar $\vec v_\infty$ cambia la magnitud de $\vec v$, así que la energía de la sonda respecto
al Sol no se conserva. Esa es la ganancia de la asistencia, y la energía sale de la órbita de
Júpiter.

En el sistema que rota con Júpiter (problema restringido de tres cuerpos) existe una
cantidad conservada, la constante de Jacobi. Lejos de Júpiter se puede aproximar por el
parámetro de Tisserand, que debería ser igual antes y después del encuentro. Pensamos usarla
para comprobar la simulación y para ver cuánto se aparta de ella el modelo de cónicas.

\section{Hipótesis iniciales}

Para el primer modelo se piensa suponer que los planetas siguen órbitas circulares y que
todo ocurre en el plano de la eclíptica. El cambio de la referencia del Sol a la de Júpiter
se haría en el borde de la esfera de influencia, y visto desde el Sol el encuentro se
trataría como instantáneo. Júpiter se toma como masa puntual y no se considera ninguna
maniobra con motor durante el sobrevuelo, de modo que la rapidez $v_\infty$ debe ser la misma
a la entrada y a la salida. Se desprecian la relatividad, la presión de radiación y los
planetas distintos de Júpiter y Saturno. Todavía falta decidir qué energía de lanzamiento
se considera disponible, y es posible que haya que probar varios valores.

\section{Dónde esperamos que falle la aproximación}

Una cuenta rápida muestra que en el borde de la esfera de influencia la atracción del Sol
sobre la sonda es unas cuatro veces mayor que la de Júpiter. Esto parece contradecir la idea
de ignorar al Sol, pero lo que importa en el sistema de Júpiter es la diferencia entre la
atracción del Sol sobre la sonda y sobre Júpiter, que es mucho menor. Aun así esperamos que
el error crezca cuando $v_\infty$ es pequeña, porque la sonda pasa más tiempo cerca del borde,
y que el resultado dependa de dónde se haga el empalme.

\section{Cómo pensamos abordarlo}

Primero se implementará el modelo de cónicas: problema de Lambert para los tramos
Tierra--Júpiter y Júpiter--Saturno, y la hipérbola para el encuentro. Con eso se buscarán
fechas de salida que permitan el viaje sin motor. Como referencia se usará la misión
Voyager~2, que salió en agosto de 1977 y pasó por Júpiter en julio de 1979, pero sus datos
solo se usarán para comparar al final. Después se integrarán numéricamente las ecuaciones
de movimiento con REBOUND, partiendo del mismo encuentro, y se compararán los dos modelos
variando $v_\infty$ y el radio de empalme. Para validar se revisará la conservación de la
energía total y de la constante de Jacobi, los casos límite (masa de Júpiter tendiendo a
cero) y la convergencia numérica.

\end{document}
